\begin{enumerate}
\item \textbf{Quantum dice and coins.}
You are given two unusual three-sided dice
which, when rolled, show either one, two, or
three spots. There are three games played with
these dice: Distinguishable, Bosons, and Fermions.
In each turn in these games, the player
rolls the two dice, starting over if required by the
rules, until a legal combination occurs. In Distinguishable,
all rolls are legal. In Bosons, a roll
is legal only if the second of the two dice shows
a number that is is larger or equal to that of the
first of the two dice. In Fermions, a roll is legal
only if the second number is strictly larger than
the preceding number. See Fig. 1.4 for a table of
possibilities after rolling two dice.
Our dice rules are the same ones that govern
the quantum statistics of noninteracting identical
particles.

\begin{enumerate}
\item Presume the dice are fair: each of the three numbers of dots shows up 1/3 of the time. For a legal turn rolling a die twice in the three games (Distinguishable, Bosons, and Fermions), what is the probability $p(5)$ of rolling a 5?




In this exercise, the quantum dice mimic distinguishable particles, bosons and fermions. Each die corresponds to a quantum particle and each face value corresponds to a quantum state. Rolling both dice corresponds to a system of two particles in state $a$ and $b$ each. We define a two particle state:

$$
|a,b> = |a>_1 \otimes |b>_2
$$

For distinguishable particles, particle 1 being in state $a$ and particle 2 being in state $b$ is different than particle 1 being in state $b$ and particle 2 being in state $a$. So $|a,b>$ and $|b,a>$ are two different two-particle states. 

For our \textit{Distinguishable} game in our quantum dice, the probability $p(5)$ of rolling a 5 is $\frac{2}{9}$ since $(2,3)$ and $(3,2)$ are two different two-particle states.

In our $\textit{Bosons}$ game, $(2,3)$ and $(3,2)$ represent the same state. By following the convention of $b \geq a$ we avoid double counting. Since we only have six possible options, $p(5) = \frac{1}{6}$

In quantum notation,
$$
|a,b> = |b,a>
$$
In the $\textit{Fermions}$ game, from the antisymmetric property of the two-particle wave function,
$$
|a,b> = -|b,a>
$$

Thus, 
$$
|a,a> = -|a,a>
$$

and $|a,a> = 0 $. Two fermions cannot occupy the same quantum state. In our \textit{Fermions} game, this is represented as the strict $b>a$ condition. Since we only have three possible options, $p(5) = \frac{1}{3}$.


\item For a legal turn in the three games, what is
the probability of rolling a double? (Hint: There
is a Pauli exclusion principle: when playing Fermions,
no two dice can have the same number
of dots showing.) Electrons are fermions; no
two noninteracting electrons can be in the same
quantum state. Bosons are gregarious (Exercise
7.9); noninteracting bosons have a larger
likelihood of being in the same state.

For the \textit{Distinguishable} game every possibility is legal. Since we have nine options and three possible doubles $p(double) = \frac{1}{3}$.


For the \textit{Bosons} game, double states are valid but we have only six possibilities. $p(double) = \frac{1}{2}$.

In the \textit{Fermions} game, the possibility of doubles is explicitly excluded so $p(double) = 0$.




\item Let us decrease the number of sides on our dice to $N=2$, making them quantum coins, with a head $H$ and a tail $T$. Let us increase the total number of coins to a large number $M$; we flip a line of $M$ coins all at the same time, repeating until a legal sequence occurs. In the rules for legal flips of quantum coins, let us make $T < H$. A legal Boson sequence, for example, is then a pattern $TTTT...HHHH...$ of length $M$; all legal sequences have the same probability. What is the probability in each of the games, of getting all the M flips of our quantum coin the same (all heads $HHHH...$ or all tails $TTTT...$ )? (Hint: How many legal sequences are there for the three games? How many of these are all the same value?)


For the \textit{Distinguishable} game, there are $2^M$ possible sequences of $H$ and $T$. Only two of these sequences are all the same, so $p = \frac{2}{2^M} = 2^{-M}$.

In the \textit{Bosons} game, there are $M+1$ legal sequences, since we can count the two sequences with all $H$ or $T$, plus the $M-1$ sequences where we have one $T$ and $M-1$ $H$'s, two $T$'s and $M-2$ $H$'s, etc. Only two of these sequences are all the same, so $p = \frac{2}{M+1}$.

In the \textit{Fermions} game, no two coins can be the same, so there are no legal sequences. Thus, $p = 0$.

The probability of finding a particular legal sequence in Bosons is larger by a constant factor due to discarding the illegal sequences. This factor is just one over the probability of a given toss of the coins being legal, $Z = \sum_{\alpha} p_\alpha$ summed over legal sequences $\alpha$. For part (c), all sequences have equal probabilities $p_{\alpha} = 2^{-M}$, so $Z_{Dist} = (2^M)(2^{-M}) = 1$, and $Z_{Boson}$ is $2^{-M}$ times the number of legal sequences. So for
part (c), the probability to get all heads or all tails is $(p_{TTT}...+p_{HHH}...)/Z$. The normalization constant $Z$ in statistical mechanics is called the partition function, and will be amazingly useful (see Chapter 6).



\item Let us now consider a biased coin, with probability $p = 1/3$ of landing H and thus $1-p = 2/3$ of landing T. Note that if two sequences are legal in both Bosons and Distinguishable, their relative probability is the same in both games. What is the probability $p_{TTT}...$ that a given toss of M coins has all tails (before we throw out the illegal ones for our game)? What is $Z_{Dist}$? What is the probability that a toss in Distinguishable is all tails? If $Z_{Bosons}$ is the probability that a toss is legal in Bosons, write the probability that a legal toss is all tails in terms of $Z_{Bosons}$. Write the probability $p_{TTT...HHH}$ that a toss has $M-m$ tails followed by $m$ heads (before
throwing out the illegal ones). Sum these
to find $Z_{Bososn}$. As M gets large, what is the probability in Bosons that all coins flip tails?


Counting both legal and illegal sequences, there are $2^M$ possible sequences of H and T.
The probability of obtaining one of those sequences (all tails) according to the given probabilities can be thought of as a Bernoulli trial with $M$ independent flips, each with probability $p = 2/3$ of being T and $1-p = 1/3$ of being H. Thus, the probability of getting all tails is:
$$
p_{TTT...} = \binom{M}{0} p^0 (1-p)^{M} = (2/3)^{M} = (2/3)^M
$$

All sequences are legal in \textit{Distinguishable}. $Z_{Dist}$ must equal 1. Each sequence has $k$ heads and $M-k$ tails. The probability of all sequences with $k$ heads and $M-k$ tails is $p_{\alpha} = \binom{M}{k} p^{k} (1-p)^{M-k}$. The sum of all probabilities is

$$
\sum_{\alpha} p_{\alpha} = \sum^{M} \binom{M}{k} p^{k} (1-p)^{M-k}
$$

From the binomial theorem,

$$
\sum_{\alpha} p_{\alpha} = \sum^{M} \binom{M}{k} p^{k} (1-p)^{M-k} = (1)^M = 1
$$

The probability that a toss in \textit{Distinguishable} is all tails is $\frac{p_{TTT...}}{Z_{Disti} = 1} = (2/3)^M$.

The probability that a toss is legal in \textit{Bosons} is $\frac{p_{TTT...}}{Z_{Bosons}}$. The probability $p_{TTT...HHH}$ that a toss has $M-m$ tails followed by $m$ heads is

$$
p_{TTT...HHH} =  p^{m} (1-p)^{M-m}
$$

With no binomial coefficient since there is only one way to have $M-m$ tails followed by $m$ heads. We sum these to find $Z_{Bosons}$

$$
Z_{Bosons} = \sum^{M}_{m=0} p^{m} (1-p)^{M-m} = (1-p)^M \sum^{M}_{m=0} (\frac{p}{1-p})^m  = (\frac{2}{3})^M \sum^{M}_{m=0} (\frac{1}{2})^m 
$$


This is a geometric series, so we can write

$$  
Z_{Bosons} = (\frac{2}{3})^M \frac{1-(\frac{1}{2})^{M+1}}{1-\frac{1}{2}} = 2(\frac{2}{3})^M (1-(\frac{1}{2})^{M+1})
$$

Finally, the probability in \textit{Bosons} that all coins flip tails is

$$
p_{TTT...} = \frac{p_{TTT...}}{Z_{Bosons}} = \frac{(2/3)^M}{2(\frac{2}{3})^M (1-(\frac{1}{2})^{M+1})} = \frac{1}{2(1-(\frac{1}{2})^{M+1})}
$$




As $M \to \infty$, $(\frac{1}{2})^M \to 0$. The probability $p_{TTT...}$ is thus

$$
p_{TTT...} = \frac{1}{2(1-0)} = \frac{1}{2}
$$
\end{enumerate}


\end{enumerate}
